\author{Brett Reynolds} %use this field for editors as well
\title{Language landscapes}
\subtitle{The ESL teacher's guide to how English works}
% \ISBNdigital{}
% \ISBNhardcover{}
% \BookDOI{}
\typesetter{Brett Reynolds}
% \proofreader{}
% \lsCoverTitleSizes{51.5pt}{17pt}
\renewcommand{\lsSeries}{tbls} 
%\newcommand{\lsSeriesTitle}{Textbooks in Language Sciences} 
%\newcommand{\lsSeriesColor}{lsYellow} 
\renewcommand{\lsSeriesNumber}{16} 
\BackBody{Language isn't a list to memorize; it's a landscape to explore. \textit{Language Landscapes} guides teachers of English across the terrain of sounds, words, sentences, and meaning, showing how patterns emerge from real use rather than from abstract rules handed down from on high. Drawing on a broad research base and the solid, functional grammar framework of \textit{The Cambridge Grammar of the English Language}, it equips future teachers to notice what speakers actually do, to test claims against evidence (including disconfirmatory evidence), and to make principled choices about pedagogy.

You'll learn to read language the way a field naturalist reads an ecosystem: attending to variation, change, register, and context. Clear examples from English (set alongside snapshots from other languages) model how to move from observation to analysis to practice. The result is a practical, intellectually honest foundation for teaching~-- one that respects learners' intuitions while sharpening teachers' judgment.

What this book helps you do:
\begin{itemize}
    \item Build a working toolkit for analyzing pronunciation, lexis, grammar, and discourse through a functional, usage-based lens.

    \item Distinguish categories from functions~-- and use that distinction to explain real learner problems clearly.

    \item Evaluate teaching materials and ``rules" by tracing the evidence behind them and looking for counter-examples.

    \item Connect linguistic insight to classroom decision-making, assessment, and feedback with confidence and care.

    \item See English in its social and global diversity, not as a single standard to be policed but as patterns to be understood.
\end{itemize}

Whether you're in pre-service training or already teaching, this book invites you to slow down, look closely, and base your instruction on what you can show. The result: explanations you can defend, feedback learners can use, and a view of English that honours variation rather than policing a single standard.}